\subsection*{一步一步建立直观}

\paragraph{1. 先由读者提出允许误差}

设你希望平方函数在 $a$ 处的差商距离 $2a$ 不超过某个容差。先确定这个输出要求，再决定步长要缩到多小。极限定义中精度与邻域的先后次序，表达的正是“无论要求多严格，我都能给出足够近的范围”。



\paragraph{2. 把图形的窄带翻译成不等式}

在本例中误差等于 $|h|$，因此可取与容差一样小的步长阈值。并非所有函数都如此简单：有时要先限制在一个较小邻域，再估计函数值的上界。证明的任务是找到能保证整段邻域都合格的条件，而非只试出几个漂亮样本。



\paragraph{3. 即使切线水平，量词仍然相同}

把观察点移到零，目标斜率也变成零。此时仍需对所有足够小的非零步长给出控制，不能因为曲线在屏幕上显得平坦就结束论证。有限像素可能把小误差隐藏起来，严格定义则要求我们继续缩小任意指定的误差。
